\section*{अध्याय \ref{ch:g10:measuring-the-earth} --- पृथ्वी को मापना}

\begin{solution}{exo:g10:measuring-the-earth:1}
जहाज़ पहले पेंदे की ओर से ओझल होता है: समतल समुद्र पर वह केवल छोटा होता जाता, और पूरा
दिखता रहता। ग्रहण में चंद्रमा पर पृथ्वी की छाया हमेशा किसी वृत्त का चाप होती है: समतल
चकती किनारे से देखे जाने पर लंबी खिंची छाया डालती। उत्तर की ओर जाने पर ध्रुव तारा क्षितिज
से ऊपर उठता है: समतल पृथ्वी पर उसकी ऊँचाई लगभग वही बनी रहती।
\end{solution}

\begin{solution}{exo:g10:measuring-the-earth:2}
\omterm{def:g10:measuring-the-earth:geodesy}{भूगणित} पृथ्वी की आकृति और आकार तथा उस पर बिंदुओं की स्थितियाँ, और समय के साथ उनके
परिवर्तन, मापता है। \omterm{def:g10:measuring-the-earth:triangulation}{त्रिभुजन} त्रिभुजों की ऐसी शृंखला से बिंदुओं की स्थिति ज्ञात करता है
जिसमें सभी कोण मापे जाते हैं; ज़मीन पर केवल एक लंबाई, आधार रेखा, मापी
जाती है।
\end{solution}

\begin{solution}{exo:g10:measuring-the-earth:3}
$C = 50 \times 5000 = 250\,000$ स्टेडियन; $250\,000 \times \qty{157.2}{m} =
\qty{39300}{km}$।
\end{solution}

\begin{solution}{exo:g10:measuring-the-earth:4}
$111/60 \approx \qty{1.85}{km}$। चौथाई याम्योत्तर: $90 \times 111 \approx
\qty{10000}{km}$।
\end{solution}

\begin{solution}{exo:g10:measuring-the-earth:5}
$a - b = 6378.137 - 6356.752 = \qty{21.385}{km}$;
$f = 21.385/6378.137 \approx 0.003353 \approx 1/298$।
\end{solution}

\begin{solution}{exo:g10:measuring-the-earth:6}
दक्षिण: $\tan\theta_1 = 0.364$, $\theta_1 \approx \qty{20.0}{\degree}$। उत्तर:
$\tan\theta_2 = 0.456$, $\theta_2 \approx \qty{24.5}{\degree}$।
$\alpha = \qty{4.5}{\degree}$, इसलिए $C = (360/4.5) \times 500 = 80 \times 500 =
\qty{40000}{km}$।
\end{solution}

\begin{solution}{exo:g10:measuring-the-earth:7}
$C = 180 - 62 - 71 = \qty{47}{\degree}$।
$AC = 10.00 \sin 71^\circ/\sin 47^\circ = 10.00 \times 0.9455/0.7314 \approx
\qty{12.93}{km}$; $BC = 10.00 \sin 62^\circ/\sin 47^\circ \approx
\qty{12.07}{km}$।
\end{solution}

\begin{solution}{exo:g10:measuring-the-earth:8}
देशांतर का परिवर्तन $10 + 20 = \qty{30}{\degree}$ है।
\qty{60}{\degree} उत्तर पर एक अंश $111.3 \times 0.5 = \qty{55.65}{km}$ है, इसलिए
दूरी $30 \times 55.65 \approx \qty{1670}{km}$ है। भूमध्य रेखा पर:
$30 \times 111.3 \approx \qty{3340}{km}$, दुगुनी दूरी।
\end{solution}

\begin{solution}{exo:g10:measuring-the-earth:9}
दृष्टि रेखा समुद्र को $T$ पर छूती है: वृत्त की स्पर्श रेखा स्पर्श बिंदु पर त्रिज्या के लंबवत
होती है, इसलिए $T$ पर का कोण
समकोण है। पाइथागोरस: $(R+h)^2 = R^2 + d^2$, इसलिए $d^2 = 2Rh + h^2$; $h$
$R$ की तुलना में नगण्य है, इसलिए $d \approx \sqrt{2Rh} = \sqrt{2 \times
6\,371\,000 \times 2} \approx \qty{5050}{m}$, लगभग \qty{5}{km}।
\end{solution}

\begin{solution}{exo:g10:measuring-the-earth:10}
वक्र मध्य अक्षांशों में, लगभग \qty{30}{\degree} से
\qty{60}{\degree} तक, सबसे खड़ा है, और भूमध्य रेखा तथा ध्रुव के पास सपाट। चपटी
पृथ्वी पर सतह ध्रुवों के पास कम वक्र होती है; सतह जितनी कम वक्र हो, ऊर्ध्वाधर (और इसलिए
अक्षांश) को एक अंश घूमने के लिए उतना ही अधिक दूर चलना पड़ता है।
\end{solution}

\begin{solution}{exo:g10:measuring-the-earth:11}
$1072.2/9.646 \approx \qty{111.15}{km}$ प्रति अंश; $90 \times 111.15
\approx \qty{10004}{km}$, वास्तविक चौथाई याम्योत्तर से \qty{2}{km} अधिक।
चाप का केंद्र लगभग \qty{46}{\degree} पर है, जहाँ एक अंश भूमध्य रेखा से ध्रुव तक
के औसत अंश से थोड़ा लंबा होता है
($10\,002.0/90 \approx \qty{111.13}{km}$)।
\end{solution}

\begin{solution}{exo:g10:measuring-the-earth:12}
\qty{7.1}{\degree} के साथ: $C = (360/7.1) \times 5000 \approx 253\,500$ स्टेडियन,
$250\,000$ के बजाय: लगभग 1.4\,\% अधिक (उसी दूरी के लिए छोटा कोण
बड़ी पृथ्वी का अर्थ रखता है)। नील नदी के किनारे-किनारे पैदल चलने वालों द्वारा क़दमों से नापी गई
दूरी कोण की तुलना में बहुत कम सटीक थी।
\end{solution}

\begin{solution}{exo:g10:measuring-the-earth:13}
$ABC$ में: $C = \qty{50}{\degree}$, $AC = 6.000 \sin 75^\circ/\sin 50^\circ
\approx \qty{7.566}{km}$। $ACD$ में: $D$ पर का कोण $180 - 50 - 68 =
\qty{62}{\degree}$ है, और $AD = 7.566 \sin 68^\circ/\sin 62^\circ \approx
\qty{7.945}{km}$।
\end{solution}

\begin{solution}{exo:g10:measuring-the-earth:14}
$\cos\varphi = 1/2$ से $\varphi = \qty{60}{\degree}$ मिलता है; $\cos\varphi = 0.1$
से $\varphi \approx \qty{84.3}{\degree}$। अक्षांश $\varphi$ पर देशांतर का एक
अंश भूमध्य रेखा पर के अंश को $\cos\varphi$ से गुणा करने पर मिलता है:
अनुपात केवल कोण पर निर्भर करता है, त्रिज्या पर नहीं।
\end{solution}

\begin{solution}{exo:g10:measuring-the-earth:15}
भूमध्य रेखा पर: $111.195 - 110.574 = \qty{0.621}{km} = \qty{621}{m}$; ध्रुव
पर: $111.694 - 111.195 = \qty{0.499}{km} = \qty{499}{m}$। प्रति अंश कई सौ मीटर की
त्रुटियाँ, जो अंश-दर-अंश जुड़ती जाती हैं, एक मीटर से कहीं बड़ी हैं: अक्षांशों और देशांतरों
से बदली गई स्थितियों के लिए दीर्घवृत्तज का ही उपयोग करना होगा।
\end{solution}

\begin{solution}{pb:g10:measuring-the-earth:1}
\textbf{1.} $31.20 - 24.09 = \qty{7.11}{\degree}$।

\textbf{2.} $7.11/360 \approx 0.0198 = 1/50.6$, एरातोस्थनीज़ के $1/50$ के
बहुत निकट।

\textbf{3.} $C = 788 \times 360/7.11 \approx \qty{39900}{km}$।

\textbf{4.} अक्षांश रेखाओं के बीच का याम्योत्तर चाप: दोपहर के कोणों का अंतर अक्षांश का
अंतर है, जो दोनों स्थानों की केवल उत्तर--दक्षिण दूरी मापता है; दूरी का पूर्व--पश्चिम
भाग कोई भूमिका नहीं निभाता।

\textbf{5.} स्थानीय दोपहर को सूर्य याम्योत्तर के तल में होता है, जो दोनों स्थानों के लिए एक
ही है; उसी दिन तारों के बीच उसकी स्थिति भी वही होती है, ताकि दोनों कोण केवल पृथ्वी की
वक्रता के कारण भिन्न हों।

\textbf{6.} $180 - 58 - 66 = \qty{56}{\degree}$।

\textbf{7.} $AC = 12.000 \sin 66^\circ/\sin 56^\circ = 12.000 \times
0.9135/0.8290 \approx \qty{13.223}{km}$।

\textbf{8.} $BC = 12.000 \sin 58^\circ/\sin 56^\circ \approx \qty{12.275}{km}$।

\textbf{9.} $D$ पर का कोण $180 - 47 - 75 = \qty{58}{\degree}$ है;
$AD = 13.223 \sin 75^\circ/\sin 58^\circ \approx \qty{15.061}{km}$।

\textbf{10.} $15.061 \cos 8^\circ \approx \qty{14.914}{km}$।

\textbf{11.} $14.914/111.13 \approx \qty{0.1342}{\degree}$, यानी
$0.1342 \times 60 \approx 8.05$ चाप-मिनट।

\textbf{12.} $111.51 - 110.57 = \qty{0.94}{km}$: \qty{940}{m}।

\textbf{13.} ध्रुव के पास दूरी के साथ ऊर्ध्वाधर धीमे घूमता है: वहाँ सतह कम वक्र है।
पृथ्वी ध्रुवों पर चपटी है और भूमध्य रेखा पर उभरी
हुई।

\textbf{14.} $a - b = \qty{21.385}{km}$; $f = 21.385/6378.137 \approx 1/298$।

\textbf{15.} साहुल दीर्घवृत्तज की सतह के लंबवत होता है। दीर्घवृत्त पर वक्र का लंब केवल
भूमध्य रेखा पर और ध्रुवों पर केंद्र से होकर जाता है; \qty{45}{\degree} पर वह
केंद्र के नीचे अक्ष को काटता है।

\textbf{16.} $10\,000\,000$ मीटर, यानी \qty{10000}{km}।

\textbf{17.} $10\,001\,966/10\,000\,000 \approx \qty{1.0001966}{m}$।

\textbf{18.} बहुत छोटा: वास्तविक करोड़वाँ भाग 1799 के मीटर का $1.0001966$ गुना है,
इसलिए 1799 का मीटर अपनी लंबाई के लगभग $1.97 \times 10^{-4}$ भाग से, यानी
0.02\,\% से, छोटा है।

\textbf{19.} बहुत छोटा। छड़ 1790 के दशक में परिकलित चौथाई याम्योत्तर का एक
करोड़वाँ भाग है; चूँकि छड़ वास्तविक चौथाई याम्योत्तर के एक करोड़वें भाग से छोटी है,
परिकलित चौथाई याम्योत्तर वास्तविक से छोटा था।

\textbf{20.} $1.0001966 - 1 = 0.0001966$ मीटर: पहला मीटर लगभग
\textbf{\qty{0.2}{mm} छोटा} था --- एक मीटर पर कागज़ की लगभग दो शीटें, दूरबीनों और
कोणों से पूरे देश में किए गए सर्वेक्षण के लिए आश्चर्यजनक सटीकता।
\end{solution}
